%% kapa: exam=1 round=37 year=2026 period=1 form=A subject=real_estate subject_ko=부동산학원론 seq=6 number=86
%% source: Q-Net 감정평가사 기출문제 (cst00309) artlSeq=5258924 file=raw/1008/5258924/2026년 제37회 감정평가사 1차 1교시.pdf
%% conversion: kordoc 4.24.51 -> tools/gen1.py (auto); review=proofread-sonnet; flags=none
\begin{question}{86}
감정평가사 A는 표준지공시지가의 감정평가를 의뢰받고 현장조사를 통해 다음과 같이 토지특성을 확인하였다. 표준지조사평가보고서상 토지특성 기재방법의 연결이 옳은 것은?
\begin{framedtext}
ㄱ. 도로접면: 폭 15 m 도로에 한면이 접하면서 자동차 통행이 가능한 폭 4 m 도로에 한면이 접하고 있는 토지 \par
ㄴ. 지형지세(고저): 간선도로 또는 주위의 지형지세보다 높고 경사도가 20°인 지대의 토지 \par
\end{framedtext}
\begin{choices}
  \item ㄱ: 중로각지, ㄴ: 고지
  \item ㄱ: 중로각지, ㄴ: 급경사
  \item ㄱ: 중로각지, ㄴ: 중경사
  \item ㄱ: 소로각지, ㄴ: 급경사
  \item ㄱ: 소로각지, ㄴ: 중경사
\end{choices}
\end{question}
\answer{2}
